% discussion.tex -- Section 13: Discussion

\section{Discussion}
\label{sec:discussion-new}

This section synthesizes findings across the game-theoretic evaluation,
training pipeline, and safety transfer experiments, and discusses
implications for AI alignment research.

\subsection{Does Strategic Reasoning Transfer to Safety?}
\label{sec:transfer-discussion}

% TODO: Analyze whether game-theoretic training improves safety benchmarks
% TODO: Discuss which game domains contribute most to transfer
% TODO: Compare GRPO vs DPO transfer effectiveness

The central question of this work is whether strategic reasoning abilities
developed through game-theoretic training generalize to broader alignment
objectives. We discuss the evidence from our safety transfer evaluation
(Section~\ref{sec:safety}) and identify which aspects of game-theoretic
reasoning---cooperation, fairness calibration, exploitation resistance---most
strongly predict safety benchmark performance.

\subsection{Coalition Dynamics and Emergent Social Norms}
\label{sec:coalition-discussion}

% TODO: Analyze emergent coalition behavior
% TODO: Discuss norm formation in repeated coalition games
% TODO: Compare enforcement mode effects on cooperation

Coalition formation (Section~\ref{sec:coalition}) creates opportunities for
emergent social norms. We discuss how different enforcement modes affect the
stability of cooperative arrangements and whether agents develop
reputation-based strategies that mirror human social behavior.

\subsection{Meta-Governance as Constitutional AI}
\label{sec:governance-discussion}

% TODO: Compare governance outcomes to Constitutional AI principles
% TODO: Analyze what rules agents propose and vote for
% TODO: Discuss democratic convergence properties

The meta-governance engine (Section~\ref{sec:governance}) operationalizes
Constitutional AI principles through democratic mechanism design. We discuss
whether agents converge on fair governance structures and how the quality of
emergent constitutions compares to hand-crafted constitutional principles.

\subsection{Open-Ended Play and Capability Evaluation}
\label{sec:openended-discussion}

% TODO: Analyze games created by agents
% TODO: Discuss implications for capability evaluation
% TODO: Address dual-use concerns

Dynamic game creation (Section~\ref{sec:dynamic}) enables novel capability
evaluation. We discuss what the games agents create reveal about their
strategic understanding and whether open-ended play provides a more robust
evaluation methodology than fixed benchmarks.

\subsection{Limitations}
\label{sec:limitations}

Several limitations should guide interpretation of our results:

\begin{itemize}[nosep]
    \item \textbf{Game abstraction}: formal games simplify real-world
          strategic interactions; transfer to naturalistic settings
          remains an open question.
    \item \textbf{Opponent diversity}: while the strategy library is
          extensive, it does not capture the full spectrum of adversarial
          or cooperative behaviors.
    \item \textbf{Scale}: current experiments use models of moderate size;
          scaling behavior is unknown.
    \item \textbf{Evaluation scope}: external benchmarks test specific
          safety dimensions; other aspects of alignment (e.g., corrigibility,
          power-seeking) are not directly measured.
    \item \textbf{Causal attribution}: improvements on safety benchmarks
          may reflect general capability gains rather than alignment-specific
          transfer.
\end{itemize}

% TODO: Expand limitations with experimental evidence
